\documentclass[12pt,a4paper]{article}

\usepackage[spanish,es-nodecimaldot]{babel}
\usepackage[utf8]{inputenc}
\usepackage[T1]{fontenc}
\usepackage{lmodern}
\usepackage{amsmath,amssymb}
\usepackage{graphicx}
\usepackage{booktabs}
\usepackage{hyperref}
\usepackage{enumitem}
\usepackage{geometry}
\usepackage{listings}
\usepackage{xcolor}

\geometry{margin=2.5cm}

\lstset{
  basicstyle=\footnotesize\ttfamily,
  breaklines=true,
  frame=single,
}

\title{Entorno RL, algoritmo y entrenamiento}
\date{}

\begin{document}
\maketitle

\section{Entorno RL: qué modela exactamente}

Archivo: \verb|src/rl_defender_env.py|

\begin{itemize}
  \item observación: \(X[i]\)
  \item acción: \texttt{0=PERMIT}, \texttt{1=BLOCK}
  \item recompensa según etiqueta real y acción
\end{itemize}

Es una formulación de decisión secuencial simple sobre muestras etiquetadas (muy cercana a contextual
bandit con coste asimétrico).

\section{Recompensa actual por defecto}

En código mantenido (\verb|train_rl_defender.py|, \verb|validate_checks.py|,
\verb|validate_leave_one_csv_out.py|):

\begin{lstlisting}[language=Python]
{
  "tp": 1.5,
  "fp": -1.5,
  "fn": -5.0,
  "omission": 0.0,
}
\end{lstlisting}

Interpretación:

\begin{itemize}
  \item bloquear ataque recompensa moderadamente
  \item bloquear benigno penaliza
  \item permitir ataque penaliza fuertemente
  \item permitir benigno queda neutro (\texttt{omission})
\end{itemize}

\section{Sobre el algoritmo principal}

Archivo: \verb|src/train_rl_defender.py|

\begin{itemize}
  \item se usa \texttt{QRDQN} (\texttt{sb3\_contrib})
  \item política: \texttt{MlpPolicy}
  \item arquitectura de red en entrenamiento principal: \texttt{[512,~256]}
  \item entorno envuelto con \texttt{DummyVecEnv} y \texttt{Monitor}
\end{itemize}

\section{Hiperparámetros relevantes en entrenamiento principal}

\begin{itemize}
  \item \texttt{learning\_rate = 1e-4}
  \item \texttt{gamma = 0.99}
  \item \texttt{tau = 1.0}
  \item \texttt{batch\_size}: depende de preset (\texttt{512} fast, \texttt{2048} full)
  \item \texttt{gradient\_steps}: depende de preset (\texttt{10} fast, \texttt{20} full)
\end{itemize}

\section{Presets y timesteps}

Código efectivo:

\begin{itemize}
  \item \textbf{fast}: \texttt{25\_000} timesteps por defecto
  \item \textbf{full}: \texttt{100\_000} timesteps por defecto
\end{itemize}

Nota: el texto de ayuda de CLI sobre timesteps puede quedar desactualizado frente al valor real
aplicado en \texttt{main()}.

\section{Artefactos que deja cada entrenamiento}

\begin{itemize}
  \item modelo: \verb|models/<RUN_ID>.zip|
  \item en \verb|runs/cicids2017/<RUN_ID>/|:
    \begin{itemize}
      \item \texttt{config.json}
      \item \texttt{metrics.json}
      \item \texttt{scaler.joblib}
      \item \texttt{train\_percentiles.npz}
    \end{itemize}
\end{itemize}

Esto permite reproducibilidad y paso consistente a Phase~2.

\section{Baseline y tuning (complemento)}

\subsection{Baseline Random Forest}

Archivo: \verb|src/baseline_random_forest.py|

Sirve para comparar contra un modelo supervisado clásico.

\subsection{Tuning con Optuna}

Archivo: \verb|src/tune_hparams.py|

Explora hiperparámetros de QRDQN y optimiza F1 de ataque.

\end{document}
